\section{分布式文件系统：目标与挑战}

文件系统基本功能：
\begin{itemize}
    \item 组织、存储、提取、管理内容+属性
    \item 命名、共享、保护文件 （类型、创建时间、权限等）
\end{itemize}

分布式场景核心需求：
\begin{itemize}
    \item 透明性 (Transparency)：访问远程文件如同本地
    \begin{itemize}
        \item 位置透明性 (Location Transparency)
        \item 位置无关性 (Location Independence)：名字解析由目录结构实现（路径名）
    \end{itemize}
\end{itemize}

核心挑战：
\begin{itemize}
    \item 网络延迟与故障
    \item 并发访问一致性
    \item 缓存协调
    \item 容错与状态管理
\end{itemize}

\section{1 抽象模型与设计问题}

\subsection{分布式文件系统抽象模型}

\subsubsection{1. 客户端}
\begin{itemize}
    \item 发起文件操作请求
    \item 可能进行本地缓存
\end{itemize}

\subsubsection{2. 服务器端}
\begin{itemize}
    \item 目录服务：路径名 + UFID 映射
    \item 文件服务：操作文件内容与属性
    \item 客户请求 → UFID (Unique File Identifier) → 服务器
\end{itemize}

\subsubsection{3. 通信}
\begin{itemize}
    \item 网络与基于 RPC
    \item 客户端与服务器
\end{itemize}

\textbf{思考}：UFID 应包含哪些信息？为何不直接用路径名作标识？(提示：防 rename race、支持 snapshot \& migration)

\subsection{分布式文件系统设计问题 (1/5)}

\subsubsection{1. 文件使用模式（经验、数据驱动设计）}
\begin{itemize}
    \item 多数文件长度 < 10KB
    \item 寿命短（临时文件）
    \item 多数文件性全本地暂存 + 普通数据
    \item 能被文件可共享 → 需一致性保障
\end{itemize}

\subsubsection{2. 命名与名字解析}
\begin{itemize}
    \item 构造名字空间（矩形、图、全局？局部？）
    \item 解析机制：路径遍历 vs. 缓存 vs. 委托
\end{itemize}

\subsection{分布式文件系统设计问题 (2/5)}

\subsubsection{3. 结构与接口}
\begin{itemize}
    \item 远程访问模型 (Remote Access)：操作粒度为 block/byte (NFS)
    \item 上载/下载 (Upload/Download)：下载/上传整个文件 (早期)
\end{itemize}

\subsection{分布式文件系统设计问题 (3/5)}

\subsubsection{4. 缓存策略与一致性}

客户端缓存一致性是指一套管理协议或规则，目的是在客户端本地访问速度与全局数据一致性之间寻求平衡。如果客户端总是直接读/写服务器，性能会很差；如果客户端只读/写本地缓存，又会出现数据不一致。

客户端缓存一致性策略：
\begin{itemize}
    \item Write-Through（直写）
    \item Delayed Write（延迟写）
    \item Write-on-Close（关闭写） (NFS 默认)
    \item Centralized Control（集中控制）
\end{itemize}

\subsubsection{Write-Through（直写）}
机制描述：客户端修改缓存数据的同时，直接同步到服务器端。
核心特点：
\begin{itemize}
    \item 数据一致性极高：服务器始终保持最新状态。
    \item 缺点：延迟较高，每次写操作都严重受限于网络带宽与延迟，性能开销大。
    \item 适用场景：对数据安全性要求极高的场景。
\end{itemize}

\subsubsection{Delayed Write（延迟写）}
机制描述：仅更新本地缓存，服务器端的同步操作被推迟到稍后时刻进行（如定时或缓冲区满）。
核心特点：
\begin{itemize}
    \item 性能极高：大部分写操作在本地内存完成。
    \item 一致性风险：若客户端在刷新到磁盘前崩溃，数据将丢失。
    \item 缺点：存在不可靠性，不适用于关键业务数据。
    \item 适用场景：对性能要求极高，且允许少量数据丢失的临时环境。
\end{itemize}

\subsubsection{Write-on-Close（关闭写）}
机制描述（NFS 默认策略）：文件打开期间的所有修改仅在本地，只有在执行 close() 操作时才将所有修改刷回服务器。
核心优势：
\begin{itemize}
    \item 平衡了性能与一致性：减少了频繁的网络交互。
    \item 符合应用习惯：符合“打开-编辑-保存”的常规流程。
\end{itemize}
潜在局限：
\begin{itemize}
    \item 若多客户端并发修改，在关闭前可能产生数据冲突或读到过期数据。
\end{itemize}

\subsubsection{Centralized Control（集中控制）}
机制描述：由服务器进行统一管理（如租约 Lease 或回调 Callback 机制），客户端遵循服务器指令。
核心特点：
\begin{itemize}
    \item 提供强一致性：支持原子操作。
    \item 动态管理：服务器可随时通知客户端作废旧缓存。
\end{itemize}
挑战：
\begin{itemize}
    \item 实现复杂，服务器开销大（需维护客户端状态）。
    \item 适用场景：共享在线系统、分布式数据库。
\end{itemize}

\subsection{分布式文件系统设计问题 (4/5)}

\subsubsection{5. 文件共享语义}
定义了多个客户端同时访问和修改时，如何处理读/写操作以保证“预期”的数据行为。
\begin{itemize}
    \item UNIX 语义：顺序一致性 (Sequential Consistency)
    \item 会话语义 (Session Semantics)：修改仅本地会话内可见。
    \item 事务语义：原子性保证。
\end{itemize}

\subsubsection{UNIX 语义 (Sequential Consistency)}
核心定义：基于本地文件系统的标准模型，提供顺序一致性。
机制：写操作立即生效，对所有用户可见。
行为：任何进程执行 write() 后，随后任何进程执行 read() 均能立即读到最新修改。
优点：编程模型简单，符合直觉。
缺点：分布式环境下性能极低（频繁网络通信）。

\subsubsection{会话语义 (Session Semantics)}
核心定义：一种折中方案，修改仅在当前会话内。
机制：
\begin{itemize}
    \item 修改仅在本地缓存进行。
    \item 关闭文件 (close) 时修改才同步回服务器。
    \item 打开文件 (open) 时才能获取最新更新。
\end{itemize}
典型应用：NFS (Network File System)。
优点：大幅减少网络交互，性能优异。
问题：存在“最后写者胜出”冲突。

\subsubsection{事务语义 (Transaction Semantics)}
核心定义：基于 ACID 原则的强一致性模型。
机制：所有操作视为一个原子事务。
特性：
\begin{itemize}
    \item 原子性：要么全成功，要么全失败。
    \item 中间状态不可见：直到 commit 提交为止。
\end{itemize}
技术支撑：通常依赖锁机制 (Locking) 或版本控制。
场景：分布式数据库、协同编辑系统。

\subsubsection{总结与对比}
\begin{center}
\begin{tabular}{lcc}
\toprule
语义 & 可见性/一致性 & 性能 \\
\midrule
UNIX & 立即 & 低 \\
会话 & 关闭时 & 高 \\
事务 & 提交时 & 低/中 \\
\bottomrule
\end{tabular}
\end{center}
选择原则：
\begin{itemize}
    \item 追求简单与一致 → UNIX 语义
    \item 追求高吞吐性能与 → 会话语义
    \item 追求强一致性和有严格要求 → 事务语义
\end{itemize}

\subsection{分布式文件系统设计问题 (5/5)}

\subsubsection{6. 断开操作与容错}
\begin{itemize}
    \item 有状态 vs. 无状态服务
    \item 服务器多副本 (Replication)
    \item 断开操作 (Disconnected Operation)
\end{itemize}

\subsubsection{服务模式：有状态 vs 无状态}

\textbf{无状态服务 (Stateless)}：服务器不保存客户端的信息（如打开的文件、锁）。\\
优势：容错简单，服务器崩溃重启后无需恢复状态。\\
挑战：每次请求都需包含完整上下文，可能增加开销。

\textbf{有状态服务 (Stateful)}：服务器记录客户端行为（如租约、打开的文件描述符）。\\
优势：性能更优，减少通信量。\\
挑战：故障恢复复杂，需精密的“状态重建”机制。

\subsubsection{服务器多副本 (Replication)}
核心目标：高可用性 (High Availability)：通过在不同节点存储多份数据，防止单点故障导致的数据丢失或服务不可用。

同步策略：
\begin{itemize}
    \item 强一致性：写操作需同步至所有副本（如 Paxos, Raft 协议）。
    \item 最终一致性：允许副本间短时延迟，通过后台异步同步。
\end{itemize}

技术挑战：
\begin{itemize}
    \item 副本间的数据不一致（Split-Brain 脑裂）。
    \item 写入性能下降（需等待多数派写入成功）。
\end{itemize}

\subsubsection{断开操作 (Disconnected Operation)}
定义：指客户端在网络不可用（或主动断开）的情况下，仍能使用读取/写访问。

关键技术：
\begin{itemize}
    \item Hoarding (预存)：在联网时先将可能使用的文件缓存到本地。
    \item Local Logging：断网期间所有修改记录在本地日志中。
    \item 重连挑战 (Reintegration)：
    \begin{itemize}
        \item 如何处理断网期间的修改冲突？
        \item 需通过版本控制或语义化合并来解决数据冲突。
    \end{itemize}
\end{itemize}

\subsubsection{权衡总结}
\begin{center}
\begin{tabular}{ll}
\toprule
设计重点维度 & 权衡 \\
\midrule
状态化 & 性能 vs 容错难度 \\
副本 & 可用性 vs 一致性成本 \\
断开 & 灵活性 vs 冲突处理复杂 \\
\bottomrule
\end{tabular}
\end{center}

原则：
\begin{itemize}
    \item 容错性往往以牺牲一部分写入性能为代价。
    \item 复杂的断开操作支持（如 Coda 文件系统）通常为已提升可用性。
\end{itemize}

\section{三、网络文件系统 NFS}

\begin{itemize}
    \item Sun 开发，事实标准
    \item POSIX 兼容
    \item 基于远程文件访问模型
    \item 通信：RPC
    \item 接口标准化：VFS (Virtual File System)
    \begin{itemize}
        \item 隐藏底层差异
        \item 本地操作 → 本地 FS
        \item 远程操作 → NFS
        \item VFS 接口上的操作或 client stub 下 server stub 传送到本地文件系统，或传送到 NFS 客户端的桩上（负责处理对存储在远程服务器上文件的访问，该访问通过 RPC 实现的）。
        \item 用户桩发出 RPC 调用
    \end{itemize}
\end{itemize}

\subsection{系统关键对象}
\textbf{文件句柄 (File Handle)}：服务器分配的不透明标识符。\\
客户端：解析路径 → 获取句柄 → 后续操作均使用句柄。\\
支持硬链接、符号链接。

\subsection{VFS}
% VFS 内容占位

\subsection{NFSv3 vs NFSv4}
\begin{center}
\begin{tabular}{lcc}
\toprule
操作 & NFSv3 & NFSv4 \\
\midrule
Create & $\checkmark$ & $\checkmark$ \\
Link & $\checkmark$ & $\checkmark$ \\
Symlink & $\checkmark$ & $\checkmark$ \\
Mkdir & $\checkmark$ & $\checkmark$ \\
Mknod & $\checkmark$ & $\checkmark$ \\
Rename & $\checkmark$ & $\checkmark$ \\
Remove & $\checkmark$ & $\checkmark$ \\
Rmdir & $\checkmark$ & $\checkmark$ \\
Open & $\times$ & $\checkmark$ \\
Close & $\times$ & $\checkmark$ \\
Lookup & $\checkmark$ & $\checkmark$ \\
\bottomrule
\end{tabular}
\end{center}

\subsection{文件句柄演进}

\textbf{持久文件句柄 (Persistent FileHandle)}：
\begin{itemize}
    \item 生命期内固定
    \item 对象删除/FS 卸载后失效
\end{itemize}

\textbf{挥发性句柄 (Volatile Handle, NFSv4)}：
\begin{itemize}
    \item 可随时失效（服务器可控）
    \item boot-time 结构：[vol=1|slot] slot：句柄表索引；
\end{itemize}

\subsection{NFS 命名与安装}

\subsubsection{1. 名字空间}
\begin{itemize}
    \item 服务器“导出”(export)
    \item 客户端“挂载”(mount) 到本地挂接点 (mount point)
    \item 可只挂载子树（非整个文件系统）
\end{itemize}

\subsubsection{安装协议}
远程文件安装协议：选定的挂接点。\\
服务器提供一组过程，由客户调用在本接点实现远程文件安装。

\subsection{NFS 文件属性}
可以用数据结构、布尔值、无符号整数或其他表示：
\begin{itemize}
    \item 强制属性 (Mandatory Attribute) --- 文件系统实现时必须支持
    \item 推荐属性 (Recommended Attribute)
    \item 命名属性 (Named Attribute)
\end{itemize}

\subsection{NFS 复合调用 (Compound Calls, NFSv4)}
将一系列操作合并为一个 RPC 调用以减少延迟。\\
(如：LOOKUP, OPEN, READ, WRITE, GETATTR 等)

属性列表：
\begin{itemize}
    \item object type：对象类型
    \item filehandle expiration：文件句柄过期时间
    \item change indicator：修改指示
    \item size：大小
    \item unix link support：UNIX 硬链接支持
    \item unix symlink support：UNIX 符号链接支持
    \item named attr：命名属性
    \item fsid：文件系统标识符
    \item lease duration：租用期限（秒）
    \item unique handle：唯一句柄
    \item readdir error：readdir 时返回的错误
    \item filehandle：文件句柄
    \item [generation]：防过期引用
\end{itemize}

操作说明：
\begin{itemize}
    \item CREATE：创建一个常规文件
    \item CREATE (特殊)：创建一个非常规文件
    \item LINK：创建一个文件的硬连接
    \item SYMLINK：创建一个文件的符号连接
    \item MKDIR：在给定目录下创建一个子目录
    \item MKNOD：创建一个特殊文件
    \item RENAME：更改文件名
    \item REMOVE：从文件系统删除一个文件
    \item RMDIR：从一个目录中删除一个空的子目录
    \item OPEN：打开一个文件
    \item CLOSE：关闭一个文件
    \item LOOKUP：根据文件名搜索文件
\end{itemize}

\subsection{持久文件句柄的格式}
\texttt{[文件系统标识符 | 文件 inode | inode 生成号 | ...]}

\subsection{VFS 层架构}
\begin{verbatim}
应用层
  |
VFS 层
  |
客户桩 --------- 服务桩
  |              |
本地文件接口    远程 SCSI 或 FC_SAN
\end{verbatim}

\subsection{NFS 远程过程调用与安全语义}

\subsubsection{1. RPC 重传}
客户发出一个 RPC 请求，并启动计时器。如果计时器超时，客户还没有收到响应，客户便使用原先相同的 XID (事务标识符) 重传原先的请求。如果文件服务器还没有处理完请求，服务器只需简单地忽略这个重传的请求。服务器通过 XID + 响应缓存实现 at-most-once 语义。

\subsubsection{2. 安全机制}
\begin{itemize}
    \item v2/v3：UNIX auth、Kerberos、Diffie-Hellman 密钥交换
    \item v4：RPCSEC\_GSS (通用安全服务 API)
    \begin{itemize}
        \item 支持多种安全机制 (Kerberos v5, SPKM, LIPKEY)
    \end{itemize}
\end{itemize}

\subsection{NFS 文件共享与加锁}

\subsubsection{1. 文件加锁}
\begin{itemize}
    \item NFS v2/v3：依赖独立协议 NLM (Network Lock Manager)
    \item NFS v4：内建加锁操作 (LOCK, LOCKU, LOCKT, RENEW)
\end{itemize}

\subsubsection{2. 访问控制}
每项指定用户/组 + 权限位 (read/write/execute/…)

% image: 图片路径未提供，保留注释

\subsection{NFS 缓存、委托与复制 (1/2)}

\subsubsection{1. 客户端缓存一致性}
\begin{itemize}
    \item 数据可缓存 (read/write)
    \item Write-on-Close：关闭时回写
    \item 缓存重验证 (Revalidation)：
    \begin{itemize}
        \item 重开已缓存文件时；比对 mtime / ctime
        \item 过期则丢弃缓存
    \end{itemize}
\end{itemize}

% image: 图片路径未提供，保留注释

\subsection{NFS 缓存、委托与复制 (2/2)}

\subsubsection{2. 委托 (Delegation, NFSv4)}
\begin{itemize}
    \item 服务器将责任委托给客户端
    \item READ delegation：客户保证期间无人写
    \item WRITE delegation：客户保证期间无人读/写
    \item 权衡：性能提升，但服务器故障时状态恢复复杂度增加（有状态代价）
\end{itemize}

\subsubsection{3. 服务器复制与迁移}
复制：仅支持整个文件系统级复制（逻辑设备）。由于复制情况下，客户在服务器之间的切换不受服务器控制，状态信息的处理是不同的。

迁移：用于负载均衡，需迁移全部状态，所有服务器状态要从原服务器传送到新服务器。NFSv4 对文件复制只提供最低限度的支持。只有整个文件系统可以被复制（包括文件、属性、目录和数据块组成的逻辑设备）。

\medskip
汇报完毕，欢迎同学们批评指教！本次课内容教材上没有。
